\section{总结与延伸}

\subsection{讲者的核心总结}

Tatsu Hashimoto 在课程结尾强调了以下几点：

\begin{enumerate}
    \item \textbf{MoE 已成标配}：在 FLOP 约束的设置下，MoE 几乎总是优于 Dense 模型。这一结论在 2025 年已经有大量经验证据支持。
    \item \textbf{稀疏性是关键}：MoE 的核心优势来自``不需要所有参数同时工作''——这使得模型可以拥有大量参数而不按比例增加计算成本。
    \item \textbf{离散路由是核心挑战}：Top-K 路由决策的不可微性使得训练变得复杂，但启发式方法（辅助损失、bias balancing）在实践中表现良好。
    \item \textbf{架构变化有限}：DeepSeek 从 V1 到 V3，核心 MoE 架构几乎未变——好的架构在小规模验证后可以直接放大。
\end{enumerate}

\subsection{全课知识图谱}

\begin{center}
\begin{tikzpicture}[node distance=0.8cm, every node/.style={draw, rounded corners, minimum width=3.5cm, minimum height=0.7cm, font=\small}]
\node (arch) {MoE 基本架构};
\node (route) [below=of arch] {路由机制};
\node (expert) [below=of route] {Expert 设计};
\node (train) [below=of expert] {训练策略};
\node (sys) [below=of train] {系统优化};
\node (ds) [below=of sys] {DeepSeek V1-V3};
\draw[->, thick] (arch) -- (route);
\draw[->, thick] (route) -- (expert);
\draw[->, thick] (expert) -- (train);
\draw[->, thick] (train) -- (sys);
\draw[->, thick] (sys) -- (ds);
\node[draw=none, right=1cm of arch, text width=5cm, font=\scriptsize] {Sparse FFN + Router, 不增 FLOPs 增参数};
\node[draw=none, right=1cm of route, text width=5cm, font=\scriptsize] {Token Choice Top-K, 线性内积路由器};
\node[draw=none, right=1cm of expert, text width=5cm, font=\scriptsize] {细粒度 Expert + 共享 Expert};
\node[draw=none, right=1cm of train, text width=5cm, font=\scriptsize] {Aux Loss + Z-Loss + Bias Balancing};
\node[draw=none, right=1cm of sys, text width=5cm, font=\scriptsize] {Expert Parallelism, Grouped GEMM};
\node[draw=none, right=1cm of ds, text width=5cm, font=\scriptsize] {MLA + MTP + Aux-Free Balancing};
\end{tikzpicture}
\end{center}

\subsection{关键 Takeaways}

\begin{importantbox}{六条核心原则}
\begin{enumerate}
    \item \textbf{MoE 用参数换 FLOPs}：更多参数、更少计算，是当前最高效的 scaling 方式
    \item \textbf{路由器越简单越好}：线性内积 + Softmax + Top-K 已经足够，复杂路由器收益有限
    \item \textbf{细粒度 Expert 是关键创新}：将 expert 切小、数量增多，是提升 MoE 性能的最有效手段
    \item \textbf{负载均衡是训练成败的关键}：没有均衡约束，Expert Collapse 几乎不可避免
    \item \textbf{好的架构在小规模就能验证}：DeepSeek V1 的架构直接放大到 V3，证明了小规模实验的价值
    \item \textbf{Expert 并不真正``专业化''}：随机路由也有效，MoE 的核心优势来自稀疏参数扩展本身
\end{enumerate}
\end{importantbox}

\subsection{拓展阅读}

\begin{itemize}
    \item Fedus et al., ``Switch Transformers: Scaling to Trillion Parameter Models with Simple and Efficient Sparsity'': \url{https://arxiv.org/abs/2101.03961}
    \item Fedus et al. 2022, ``A Review of Sparse Expert Models in Deep Learning'': \url{https://arxiv.org/abs/2209.01667}
    \item Shazeer et al. 2017, ``Outrageously Large Neural Networks: The Sparsely-Gated Mixture-of-Experts Layer'': \url{https://arxiv.org/abs/1701.06538}
    \item DeepSeek-MoE: \url{https://arxiv.org/abs/2401.06066}
    \item DeepSeek-V2: \url{https://arxiv.org/abs/2405.04434}
    \item DeepSeek-V3: \url{https://arxiv.org/abs/2412.19437}
    \item OLMoE: \url{https://arxiv.org/abs/2409.02060}
    \item Mixtral of Experts: \url{https://arxiv.org/abs/2401.04088}
    \item Zoph et al., ``ST-MoE: Designing Stable and Transferable Sparse Expert Models'': \url{https://arxiv.org/abs/2202.08906}
\end{itemize}

\end{document}
